% GENERATED FILE. Edit canonical lessons, then run npm run generate:books.
% canonical-source-hash: fnv1a64:778ca7f4db8b5a78
% canonical-lessons: ZH-C355-xiejue, ZH-C355-haoyi, ZH-C355-ken, ZH-C355-fouren, ZH-C355-dangmian

\chapter{Thanks, but No}
\label{ch:zh-thanks-but-no}
\hlchaptermodality{355}

\begin{chapteropening}
\textbf{By the end of this chapter:} \emph{I can turn something down politely, say someone won't do something, and say it to their face.}

Complete the last lesson of chapter 355: I can turn something down politely, say someone won't do something, and say it to their face.
\end{chapteropening}

\section[xièjué]{xièjué — to decline politely}

\label{lesson:ZH-C355-xiejue}



\begin{quote}
Before the new one: say the Mandarin for what something is about, then the Mandarin for to have a serious talk.
\end{quote}

\subsection*{You'll want to know: xièjué}
\emph{xièjué} — \textquotedblleft{}to decline politely\textquotedblright{}.

\begin{quote}
\emph{wǒ qǐng tā chī fàn tā xièjué le} — I invited him to dinner, but he politely declined.
\end{quote}

Its characters come later; for now it is a word for the ear.

\begin{cousinweb}
\emph{xièjué} is built from \zh{谢} and \emph{jué}, both inherited from Middle Chinese, \emph{zjaeH} and \emph{dzjwet} (Guangyun, in Baxter's transcription), which glosses the first 'to take leave' and the second 'to cut off'.
\end{cousinweb}

\begin{grammarlens}[title={What you've built}]
One more word for giving reasons.
\end{grammarlens}

\subsection*{Your turn}
\begin{itemize}
\raggedright
  \item \emph{Say it:} \emph{xièjué}
  \item \emph{Say it:} \emph{xièjué}, once more
  \item \emph{Say it:} \emph{wǒ qǐng tā chī fàn tā xièjué le} — I invited him to dinner, but he politely declined.
  \item \emph{Recall:} say the Mandarin for what something is about, then the Mandarin for to have a serious talk, then say \emph{xièjué} again
\end{itemize}

\begin{tcolorbox}[breakable,colback=teal!4,colframe=teal!35!black,title={Before you move on}]
What does xièjué mean? (\textquotedblleft{}To decline politely\textquotedblright{}.) Say it once more.
\end{tcolorbox}
\section[\texorpdfstring{\zh{好意} (hǎoyì)}{好意 (hǎoyì)}]{\zh{好意} (hǎoyì) — good intentions; a kind thought}

\label{lesson:ZH-C355-haoyi}



\begin{quote}
Before the new one: say the Mandarin for to have a serious talk, then the Mandarin for to decline politely.
\end{quote}

\subsection*{You'll want to know: \zh{好意}}
\textbf{\zh{好意}} — \emph{hǎoyì} — \textquotedblleft{}good intentions; a kind thought\textquotedblright{}.

\begin{quote}
\emph{xièxie nǐ de hǎoyì wǒ zìjǐ néng huí jiā} — Thanks for the kind thought, but I can get home on my own.
\end{quote}

\begin{cousinweb}
\textbf{\zh{好意}} (\emph{hǎoyì}) is built from \zh{好} 'good' and \zh{意} 'intent', both inherited from Middle Chinese, \emph{xawX} and \emph{'iH} (Guangyun, in Baxter's transcription), which glosses the first 'good' and the second 'will'.
\end{cousinweb}

\begin{grammarlens}[title={What you've built}]
One more word for giving reasons.
\end{grammarlens}

\subsection*{Your turn}
\begin{itemize}
\raggedright
  \item \emph{Say it:} \emph{hǎoyì}
  \item \emph{Say it:} \emph{hǎoyì}, once more
  \item \emph{Say it:} \emph{xièxie nǐ de hǎoyì wǒ zìjǐ néng huí jiā} — Thanks for the kind thought, but I can get home on my own.
  \item \emph{Recall:} say the Mandarin for to have a serious talk, then the Mandarin for to decline politely, then say \emph{hǎoyì} again
\end{itemize}

\begin{tcolorbox}[breakable,colback=teal!4,colframe=teal!35!black,title={Before you move on}]
What does \zh{好意} mean? (\textquotedblleft{}Good intentions; a kind thought\textquotedblright{}.) Say it once more.
\end{tcolorbox}
\section[kěn]{kěn — to be willing to; to agree to}

\label{lesson:ZH-C355-ken}



\begin{quote}
Before the new one: say the Mandarin for to decline politely, then the Mandarin for a kind thought.
\end{quote}

\subsection*{You'll want to know: kěn}
\emph{kěn} — \textquotedblleft{}to be willing to; to agree to\textquotedblright{}.

\begin{quote}
\emph{tā bù kěn gàosu wǒ wèishénme} — He won't tell me why.
\end{quote}

Its characters come later; for now it is a word for the ear.

\begin{cousinweb}
\emph{kěn} is inherited from Middle Chinese \emph{khongX} (Guangyun, in Baxter's transcription), which glosses it 'to permit' and, quoting an older dictionary, 'the flesh between the bones'. How it came to mean being willing is not settled.
\end{cousinweb}

\begin{grammarlens}[title={What you've built}]
One more word for giving reasons.
\end{grammarlens}

\subsection*{Your turn}
\begin{itemize}
\raggedright
  \item \emph{Say it:} \emph{kěn}
  \item \emph{Say it:} \emph{kěn}, once more
  \item \emph{Say it:} \emph{tā bù kěn gàosu wǒ wèishénme} — He won't tell me why.
  \item \emph{Recall:} say the Mandarin for to decline politely, then the Mandarin for a kind thought, then say \emph{kěn} again
\end{itemize}

\begin{tcolorbox}[breakable,colback=teal!4,colframe=teal!35!black,title={Before you move on}]
What does kěn mean? (\textquotedblleft{}To be willing to; to agree to\textquotedblright{}.) Say it once more.
\end{tcolorbox}
\section[fǒurèn]{fǒurèn — to deny}

\label{lesson:ZH-C355-fouren}



\begin{quote}
Before the new one: say the Mandarin for a kind thought, then the Mandarin for to be willing.
\end{quote}

\subsection*{You'll want to know: fǒurèn}
\emph{fǒurèn} — \textquotedblleft{}to deny\textquotedblright{}.

\begin{quote}
\emph{wǒ juéde tā bù yīnggāi fǒurèn} — I don't think he should deny it.
\end{quote}

Its characters come later; for now it is a word for the ear.

\begin{cousinweb}
\emph{fǒurèn} is built from \emph{fǒu} and \emph{rèn}, both inherited from Middle Chinese, \emph{pjuwX} and \emph{nyinH} (Guangyun, in Baxter's transcription), which glosses the first, quoting an older dictionary, 'not', and the second 'to recognize'.
\end{cousinweb}

\begin{grammarlens}[title={What you've built}]
One more word for giving reasons.
\end{grammarlens}

\subsection*{Your turn}
\begin{itemize}
\raggedright
  \item \emph{Say it:} \emph{fǒurèn}
  \item \emph{Say it:} \emph{fǒurèn}, once more
  \item \emph{Say it:} \emph{wǒ juéde tā bù yīnggāi fǒurèn} — I don't think he should deny it.
  \item \emph{Recall:} say the Mandarin for a kind thought, then the Mandarin for to be willing, then say \emph{fǒurèn} again
\end{itemize}

\begin{tcolorbox}[breakable,colback=teal!4,colframe=teal!35!black,title={Before you move on}]
What does fǒurèn mean? (\textquotedblleft{}To deny\textquotedblright{}.) Say it once more.
\end{tcolorbox}
\section[\texorpdfstring{\zh{当面} (dāngmiàn)}{当面 (dāngmiàn)}]{\zh{当面} (dāngmiàn) — to someone's face; in person}

\label{lesson:ZH-C355-dangmian}



\begin{quote}
Before the new one: say the Mandarin for to be willing, then the Mandarin for to deny.
\end{quote}

\subsection*{You'll want to know: \zh{当面}}
\textbf{\zh{当面}} — \emph{dāngmiàn} — \textquotedblleft{}to someone's face; in person\textquotedblright{}.

\begin{quote}
\emph{nǐ zuìhǎo dāngmiàn gēn tā shuō} — You'd better say it to his face.
\end{quote}

\begin{cousinweb}
\textbf{\zh{当面}} (\emph{dāngmiàn}) is built from \zh{当} and \zh{面} 'the face', both inherited from Middle Chinese, \emph{tang} and \emph{mjienH} (Guangyun, in Baxter's transcription), which glosses the first 'to match' and 'to preside', and the second 'to face' and 'in front'.
\end{cousinweb}

\begin{grammarlens}[title={What you've built}]
One more word for giving reasons.
\end{grammarlens}

\subsection*{Your turn}
\begin{itemize}
\raggedright
  \item \emph{Say it:} \emph{dāngmiàn}
  \item \emph{Say it:} \emph{dāngmiàn}, once more
  \item \emph{Say it:} \emph{nǐ zuìhǎo dāngmiàn gēn tā shuō} — You'd better say it to his face.
  \item \emph{Recall:} say the Mandarin for to be willing, then the Mandarin for to deny, then say \emph{dāngmiàn} again
\end{itemize}

\begin{tcolorbox}[breakable,colback=teal!4,colframe=teal!35!black,title={Before you move on}]
What does \zh{当面} mean? (\textquotedblleft{}To someone's face; in person\textquotedblright{}.) Say it once more.
\end{tcolorbox}
